\subsection{Continuation of row (a): domains and derived comparisons}
\label{sec:a-progress-20260907}
This continuation integrates P123.063--123. The detailed constructions
and proofs are archived in the dated update directories. The statements
below retain their specified domains; they do not assert completion of
row (a).

P123.063 extends the logarithmic comparison to $t=1$ on sufficiently
small balls: $p^a\mathbb Z_p$, with $a\geq1$ for odd $p$ and $a\geq2$
for $p=2$. P123.064--069 construct bounded twist complexes, relative
completed bars, cyclic comparisons and their multibar coherence.
P123.070--072 distinguish the relative scalar retract from finite
perfectness over $\mathbb C$: the scalar rapid algebra already has an
infinite-dimensional degree-zero shadow. P123.073--074 realize graph
kernels and their half-density adjoints, retaining the distinction
between coefficient effectivity and positive-semidefinite squares.

P123.082 computes morphisms between elementary twists and shows why
finite cones of those twists alone do not supply the mixed source.
P123.083 supplies a larger localized sector: for the stated rapid
coefficient ideals,
\[
 I_U\widehat\otimes_A I_V\simeq I_{U\cap V},
\]
with strict exactness of the corresponding augmented bar. Nevertheless,
P123.084 shows that its original continuous degree-zero shadow vanishes
on open sets avoiding the origin. P123.085--090 therefore construct
stronger integrable modules, external boundary traces, covariant normal
actions, periodization and cutoff partial isometries. The resulting
quotient representation is identified with the induced representation
of P123.043. This is a comparison on the declared boundary sector;
it does not yet glue all strata of the arithmetic source.

\subsection{Tate excess and the topology of corner transport}
P123.075--081 compute the full Tate bar excess by strict comparisons.
For $r$ prime coordinates, a character self-intersection has the
exterior-algebra Tor groups $\bigwedge^\bullet\mathbb C^r$; distinct
Tate characters give acyclic mixed fibers. The multiplication cone
retains two copies of $\bigwedge^k\mathbb C^r$ in degree $-k-1$.
The determinant transitions are explicit, but do not remove the
underlying excess complexes or provide their arithmetic metric.
The strict exactness arguments use convergent residual fillings;
no continuous linear splitting is inferred from those fillings.

P123.091 identifies the cutoff corner with $C(Y)$ but proves that the
reverse Morita tensor does not surject onto the original full rapid
algebra. P123.092--093 construct an orbit-decay ideal $J$ for which
both completed tensor identities hold:
\[
 Je\widehat\otimes_{eJe}eJ\simeq J,
 \qquad eJ\widehat\otimes_J Je\simeq eJe=C(Y).
\]
P123.094 identifies both its internal and ambient commutator quotients
with $C(Y)$ by periodization. P123.095 transports perfect complexes in
this corner sector to actual vector bundles and their graded determinant
lines. This is a determinant construction on a specified sector, rather
than a determinant of the full mixed diagonal.

\subsection{Smooth diagonals and unit-level completion}
P123.096--097 separate the moving and fixed continuous bars. Moving
translations admit explicit contractions; the fixed continuous bar has
a strict comparison with $C(Y)[0]$. P123.098 explains why these fibers
cannot form a perfect family across the fixed parameter: their virtual
ranks jump. At a finite unit level the smooth diagonal instead retains
its first-order term. Writing $B_N=C^\infty(Y_N)$, P123.099 gives
\[
 H_0=B_N,\qquad H_1=B_N\,dt,\qquad H_j=0\quad(j\geq2).
\]
Under refinement the pulled-back diagonal is the union of all deck
transformation graphs. P123.100 retains every summand in its derived
comparison; projection to the identity graph does not erase the others.

P123.101 constructs the nuclear Fourier/conductor algebra
$\mathcal S(Y)$ and its completed tensor square. P123.102 identifies
its completed first bar homology with $\mathcal S(Y)\,dt$.
P123.118, incorporated below, resolves the remaining higher homology
of this particular completed smooth bar.
P123.103 constructs the continuous inverse of the tangential derivative
on the mean-zero subspace and moving-bar contractions. It also proves
that the real-direction heat operator has infinite trace on the
associated full Hilbert space. Nuclearity of the coefficient algebra
alone therefore does not supply the missing Hilbert trace.

\subsection{Relative analytic determinants and composition}
P123.104 constructs the tangential Fredholm determinant with convergence
from finite unit levels. P123.105--108 treat the normal compression,
including both valuation directions, nonlinear bulk subtraction and
all-order trace estimates. Opposite shifts exhibit a nontrivial
composition correction $1/(1-st)$; multiplicativity of compressed
operators cannot be assumed.

P123.109 performs the coupled construction on the actual induced
representation. On each fixed external Bruhat test domain put
$B_L(F)=E_LR(F)E_L$. Its finite collection of trace-class
operator-valued Laurent symbols is denoted $P_{F,\chi}(w)$. For $z$
sufficiently small define
\[
 b_F(z)=\sum_\chi\int_0^{2\pi}
 \log\det\bigl(1+zP_{F,\chi}(e^{i\theta})\bigr)
 \frac{d\theta}{2\pi}.
\]
The all-order endpoint estimates give the locally uniform limit
\[
 \Delta_F(z)=\lim_{L\to\infty}
 \det(1+zB_L(F))\exp\bigl(-(L+1)b_F(z)\bigr),
 \qquad (\log\Delta_F)'(0)=W_q(F).
\]
It is a nonvanishing holomorphic germ and respects the prescribed
normalized scalar semilocal pullbacks. Section changes have this
invariance when the cutoffs are transported with the unitary comparison.
These assertions do not identify arbitrary independently chosen cutoffs.

P123.110 constructs a trace-class endpoint Fredholm correction
$\mathcal A(F,G)$ for the convolution-unitization product:
\[
 \Delta_{F+G+F*G}(1)
   =\Delta_F(1)\Delta_G(1)\mathcal A(F,G).
\]
Here all arguments lie in a sufficiently small fixed test neighborhood.
The correction satisfies the associativity cocycle, transpose and
prescribed semilocal identities. Its endpoint Hilbert space is a finite
sum of tangential Hilbert spaces, not generally finite-dimensional.
An analytic cocycle on this test domain is not yet the geometric
biextension required by mixed Riemann--Roch.

\subsection{Intersection, boundary energy and the missing metric comparison}
P123.111 identifies the bilinear variation of this correction with the
stabilized endpoint trace
\[
 \mathcal B(F,G)=\operatorname{Tr}
 \bigl(E_LR(F)(1-E_L)R(G)E_L\bigr).
\]
It is symmetric, and $\mathcal B(F,G^*)$ is a positive Hermitian
boundary energy. P123.112 proves that the radical of the analogous
local normal anomaly pairing is not an ideal: a unit-character
projector has zero anomaly against every test, while its product with
a valuation shift has nonzero anomaly. Thus that pairing cannot simply
be substituted for the numerical intersection form of row (a).
P123.113 computes the distinct local trace-product pairing: on the
compact test algebra its radical is zero and the form is indefinite.
Neither statement supplies a global sign theorem.

Let $\lambda(F)=\log\Delta_F(1)$ near zero. P123.114 proves the
second-variation identity
\[
 W_q(F*G)=\mathcal B(F,G)-d^2\lambda_0(F,G).
\]
Indeed, the difference between the trace of the compressed product and
the compressed convolution is exactly the endpoint defect above;
the bulk-subtracted former trace is $-d^2\lambda_0(F,G)$.
This connects the analytic determinant to the relative trace-product
pairing with its actual normalization.

P123.115 adds explicit metric data. With
$X_L(F)=(1-E_L)R(F)E_L$, the positive Fredholm determinant
\[
 h(F)=\det\bigl(1+X_L(F)^*X_L(F)\bigr)\geq1
\]
is exactly independent of sufficiently large $L$ on the fixed test
domain. It defines the metric $\|1\|_F^2=h(F)^{-1}$ on a trivial line.
Its potential has Levi form
\[
 \left.\partial_z\partial_{\bar w}\log h(zF+wG)\right|_0
       =\mathcal B(F,G^*).
\]
By contrast, $2\operatorname{Re}\lambda$ is pluriharmonic. The
holomorphic Hessian in the preceding identity is therefore not the
curvature of the relative determinant modulus. The Gram metric
respects the prescribed semilocal pullbacks, but a multiplicative
isometry law and identification with the intrinsic periodic ruling
metric remain to be constructed.

P123.116 extends this curvature computation to every amplitude. Set
$X=X_L(F)$, $A=X_L(U)$, $C=X_L(V)$, $P=1+X^*X$ and
$Q=1+XX^*$. Trace-norm differentiation gives
\[
 K_F(U,V)=\operatorname{Tr}(P^{-1}C^*Q^{-1}A),\qquad
 K_F(U,U)=\|Q^{-1/2}AP^{-1/2}\|_{\mathrm{HS}}^2.
\]
Indeed the mixed derivative of $\log\det P$ is
$\operatorname{Tr}(P^{-1}C^*[1-XP^{-1}X^*]A)$, and the bracket
is $Q^{-1}$. Both inverse square roots are boundedly invertible.
Consequently the null directions are precisely $X_L(U)=0$,
independent of the base point, and
\[
 (1+\|X_L(F)\|^2)^{-2}\|X_L(U)\|_{\mathrm{HS}}^2
 \leq K_F(U,U)\leq\|X_L(U)\|_{\mathrm{HS}}^2.
\]
The affine test-family Gram determinant stabilizes before differentiation,
so this is also an exact cutoff-independent curvature formula. It
respects the prescribed semilocal transport. Holomorphic changes of
frame add only pluriharmonic terms to the potential and cannot change
its kernel. In the local normal model of P123.112 the unit-character
projector remains a null direction at every base point, while its
product with the valuation shift remains non-null. Thus neither moving
the base point nor changing holomorphic frame makes this particular
local curvature kernel a convolution ideal. Additional geometric metric
data are still necessary for the proposed arithmetic comparison.

\subsection{The composition correction as a normalized graph metric}
P123.117 links the analytic composition correction to metric data.
For a sufficiently small test $F$, put
\[
 a_F=1+E_LR(F)E_L,\qquad c_F=(1-E_L)R(F)E_L,
 \qquad Y_F=c_Fa_F^{-1}.
\]
The represented operator $1+R(F)$ maps the cutoff space onto the graph
of $Y_F$. The exact factorization
\[
 1+E_LR(F^*\mathbin{\diamond}G)E_L
       =a_F^*(1+Y_F^*Y_G)a_G,
 \qquad F\mathbin{\diamond}G=F+G+F*G,
\]
identifies its finite-cutoff determinant correction with
$\det(1+Y_F^*Y_G)$. Consequently
\[
 \mathcal A(F^*,G)=\lim_L\det(1+Y_F^*Y_G),\qquad
 \widehat h(F)=\mathcal A(F^*,F)\geq1.
\]
Convergence follows from the endpoint resolvent estimates: if
$c_F^*c_F=W_LdW_L^*$, the diagonal determinant equals
\[
 \det\bigl(1+d^{1/2}W_L^*a_F^{-1}(a_F^*)^{-1}W_Ld^{1/2}\bigr),
\]
whose positive perturbation converges in trace norm. These assertions
hold in the common small test neighborhood, not at arbitrary amplitude.

The mixed kernel is positive definite. At finite cutoff it is the
Gram kernel of
$\Psi(Y)=(1,Y,\bigwedge^2Y,\ldots)$ in the Hilbert direct sum of
Hilbert--Schmidt exterior-map spaces, since
\[
 \langle\Psi(Y_F),\Psi(Y_G)\rangle=\det(1+Y_F^*Y_G),
 \qquad \|\Psi(Y)\|^2=\det(1+Y^*Y)<\infty.
\]
Finite-rank approximation proves these identities; positivity passes
to each finite Gram matrix in the limit. This provides an actual
analytic kernel, rather than assuming an intersection form positive.

The existing cocycle and adjoint identities now yield
\[
 \Delta_{F^*\mathbin{\diamond}F}(1)
          =|\Delta_F(1)|^2\widehat h(F).
\]
Thus $\widehat h^{-1}$ is a metric explicitly tied to the relative
determinant composition law. Its potential is plurisubharmonic on
finite-dimensional test slices, and its Levi form at zero is the
boundary energy $\mathcal B(U,V^*)$. All these identities respect
the prescribed semilocal pullbacks. They do not identify the line with
the arithmetic determinant or its curvature with the trace-product
intersection. In particular the local radical obstruction at zero
remains; positivity of this analytic kernel is not the Weil sign law.

\subsection{Continuous contraction of the higher smooth diagonal}
P123.118 closes the higher-degree question left by P123.102 for the
normalized completed Hochschild complex of $S=\mathcal S(Y)$. It
constructs continuous $S$-linear maps $H_n:C_n\to C_{n+1}$ with
\[
 bH+Hb=1-s\pi,\qquad
 C_\bullet\simeq S[0]\oplus S\,dt[1].
\]
Here $\pi_0=1$, $\pi_1(m\otimes[f])=mDf\,dt$, and
$s_1(h\,dt)=h\chi_0^{-1}\otimes[\chi_0]/(i\tau_0)$.
In particular all higher homology vanishes and all boundary images
are closed, with continuous right inverses on those images.

The proof supplies the estimates needed for completion. Separating the
total Fourier character identifies Hochschild chains with group-bar
chains of the character group and a free $S$ coefficient. At unit
level $N$, the subgroup generated by $\chi_0$ has finite index. Choose
coset representatives $s(k)$ with frequencies in $[0,\tau_0)$ and put
$r(g)=g-s(\bar g)$. Averaging homogeneous simplices
\[
 [g_0,\ldots,g_n]\longmapsto
 \frac1{|K_N|}\sum_{k\in K_N}
 [r(g_0+s(k)),\ldots,r(g_n+s(k))]
\]
gives a transfer to the infinite cyclic subgroup. An averaged prism
is a homotopy from its inclusion-composite to the identity. The cyclic
bar contracts to the cellular integer line: edges map to their oriented
unit-edge paths, and higher homotopies cone the recursively formed
cycles from their first vertex. Vertex ranges remain in the original
integer interval and coefficient sums grow at most linearly in its
diameter, with constants depending on degree.

These finite-level homotopies need not be compatible as chosen.
P123.118 therefore assembles the global homotopy recursively, starting
with the degree-one filling of P123.102. For a Fourier basis simplex
$\sigma$ in degree $n\geq2$, take its least unit level $N$ and set
\[
 z_\sigma=\sigma-H_{n-1}b\sigma,\qquad
 H_n\sigma=h_n^{(N)}z_\sigma.
\]
The residual $z_\sigma$ is a cycle in that same level; the finite-level
homotopy fills it exactly. Coset representatives, unit paths and the
recursion give, for each degree and seminorm, a finite exponent $B(n,m)$
such that
\[
 \|H_n[g_1|\cdots|g_n]\|_m
       \leq C_{n,m}\prod_i W(g_i)^{B(n,m)}.
\]
These polynomial bounds extend the maps continuously to the weighted
Fourier completion and extend the homotopy identity by density.
Thus vanishing is proved on the completed complex itself. This gives a
bounded finite free $S$-module model for this boundary shadow; it does
not assert a finite projective resolution of the diagonal over
$S\widehat\otimes S$, full smooth orbit Morita transport, or arithmetic
mixed Riemann--Roch.

\subsection{A common smooth source and the cyclic boundary comparison}
P123.119 proves that the projection of P123.118 is a strict morphism
of mixed complexes: Connes' operator maps under it to the physical
derivative $D:S\to S\,dt$. The kernel has the continuous Hochschild
contraction already constructed. For a formal $(b+uB)$-cycle in that
kernel, the recursive filling
\[
 y_j=H(x_j-By_{j-1})
\]
proves comparison in the coefficientwise formal negative and periodic
cyclic totalizations. This does not assert entire cyclic convergence.
The Green operator of P123.103 then computes periodic homology:
\[
 HP_{2j}(S)=K_0,\qquad HP_{2j+1}(S)=K_0\,dt,
 \qquad K_0=\ker D=C(U/\overline{\langle q\rangle}).
\]
The unit subgroup $\overline{\langle q\rangle}$ is open: for a suitable
positive integer $m$, put $b=v_p(q^m-1)<\infty$ in the logarithm's
small-ball domain; closure of the powers of $q^m$ contains
$1+p^b\mathbb Z_p$. Thus $K_0$ is finite-dimensional. If its dimension
is $r$, the original trace is $(\log q)/r$ times the sum on these
$r$ quotient points. The finite periodic determinant is thereby defined,
but remains distinct from a Hilbert heat determinant or arithmetic RR.

P123.120 compares this smooth algebra with the original boundary
composition. Impose all real derivative and unit-conductor seminorms
on the orbit-decay source $J$ to obtain $J_{\rm sm}$. Smooth finite-unit
cutoffs give
\[
 eJ_{\rm sm}e=S(Y),\qquad
 J_{\rm sm}e\widehat\otimes_S eJ_{\rm sm}=J_{\rm sm},\qquad
 eJ_{\rm sm}\widehat\otimes_{J_{\rm sm}}J_{\rm sm}e=S.
\]
The frame estimates of P123.093 survive differentiation and conductor
weights: translated fixed frames have fixed unit level and bounded
real derivatives, while arbitrarily high orbit weights absorb their
label growth. The inverse tensor series therefore converge in the
smooth topology. These identities provide monoidal transport on the
specified firm complete bimodules, with the original unit and transpose,
and transport the bounded corner-perfect determinant sector.

P123.121 places that boundary inside one coefficient algebra across
all finite-zero strata. For each prime $p$, define $F_p$ by absolutely
summable shell/unit Fourier coefficients with every weight
\[
 \sum_{j,\kappa}p^{l|j|}c_p(\kappa)^N|f_{j,\kappa}|,
 \qquad E_p=\mathbb C\chi_p\oplus F_p,
 \quad \chi_p=1_{\mathbb Z_p}.
\]
The functions in $F_p$ extend by zero at the origin and decay at both
valuation ends. Evaluation at zero splits $E_p\to\mathbb C$.
Dilation preserves $F_p$ and retains the actual correction
\[
 \alpha_{p^h v}\chi_p-\chi_p
       =1_{p^h\mathbb Z_p}-1_{\mathbb Z_p}\in F_p.
\]
Let $E_{\mathbb R}$ consist of smooth real functions with every
exponentially weighted derivative bounded. The common coefficients and
source are
\[
 E_P=E_{\mathbb R}\widehat\otimes\widehat\bigotimes_{p\in P}E_p,
 \qquad A_{\rm sm}(P)=\operatorname{Rapid}_{\Gamma_P}(E_P).
\]
Weighted sequence estimates prove nuclearity and completeness; dilation
has polynomial growth in the original lattice height, so the original
crossed multiplication and involution are continuous. The coefficient
realization is injective into the original continuous source.

The tensor splittings at zero give strict split source extensions
for the full zero-count filtration. Their exactly-$k$ layers are
\[
 \bigoplus_{|J|=k}\operatorname{Rapid}_{\Gamma_P}
 \left(E_{\mathbb R}\widehat\otimes
       \widehat\bigotimes_{p\notin J}F_p\right).
\]
In the two-prime single-zero case this is precisely $J_{\rm sm}$,
with equivalent seminorms. Thus the boundary comparison is a comparison
on an actual layer of this common source. Neither the generic nor the
multiple-zero layers have been discarded.

For $P\subset T$, the coefficient map
$f\mapsto f\prod_{\ell\in T\setminus P}\chi_\ell$, together with the
old-lattice inclusion, is an injective continuous algebra map. Old
lattice elements act as units at the added primes and preserve their
cutoffs; this proves multiplicativity after crossing. These maps compose
exactly on triples and extend to the coefficient unitizations with the
actual old gauge unitaries. They are coherent semilocal source maps.

The construction now joins the boundary algebra to a common stratified
source with its attaching corrections. It still requires a localized
derived correspondence category on the full square and its shadow,
cyclic/excision comparisons on that domain, the numerical and effectivity
constructions, and the mixed determinant identity with its intrinsic
ruling metric. Algebraic source gluing alone does not prove those claims.

\subsection{Derived composition on the common algebra}
P123.122 constructs the relative derived category of bounded-above
complete Fr\'echet bimodule complexes over $A=A_{\rm sm}^u(P)$.
Admissible sequences split continuously after forgetting both actions.
The free bimodules $A\widehat\otimes V\widehat\otimes A$ lift maps
through these surjections. Left and right completed bars give relatively
projective resolutions $Q(M)$: a term with internal degree $i$ and bar
degrees $r,s$ lies in degree $i-r-s$, so each total degree has only
finitely many summands. Unit insertion proves the augmentation is an
underlying continuous chain-homotopy equivalence.

Composition is
\[
 M\star N=Q(M)\widehat\otimes_A Q(N).
\]
Free bimodules are flat for the stated split exact structure; the
bounded-above totalization preserves this property by a finite-in-each-
degree contraction argument. Consequently composition is independent
of the chosen resolutions. Balanced tensor association gives the
associator, its pentagon and the two-sided unit $A$. The relative
diagonal complex
\[
 \operatorname{Tr}_A(M)
    =\operatorname{Tot}_{i,n}\bigl(M^i\widehat\otimes A^{\widehat\otimes n}\bigr),
 \qquad \deg=i-n,
\]
has continuous differentials, respects the relative replacements and
cones, and the multiple-bar rotation proves
$\operatorname{Tr}_A(M\star N)\simeq\operatorname{Tr}_A(N\star M)$.
On degreewise nuclear modules it is a degreewise nuclear complex;
no finite determinant of an arbitrary unbounded complex is inferred.

For an inclusion $A\to B$, the well-defined induction
$B\widehat\otimes_A^{\mathbf L}M\widehat\otimes_A^{\mathbf L}B$
is coherent on chart triples, but is not asserted strong monoidal.
The middle $A$ in a source composition becomes a middle $B$ after
two induced objects are composed; the resulting comparison need not
be invertible. Likewise the source operation
$a\bar m b=\overline{b^*ma^*}$ is a covariant conjugate-opposite
involution reversing tensor order. It does not supply the contravariant
duality required by A.4. These distinctions retain the actual descent
and duality obligations. Geometric localization on the full square,
the suitable determinant domain, numerical equivalence, effectivity
and the mixed RR metric identity remain required.

\subsection{An equivariant repair of the semilocal composition map}
P123.123 proves that the induction unit comparison fails on the actual
common source. For a new prime $\ell$, the continuous derivation
$\delta_\ell(\sum b_\gamma U_\gamma)=\sum v_\ell(\gamma)b_\gamma U_\gamma$
vanishes on the old algebra $A$. Hence
$b\otimes c\mapsto b\delta_\ell(c)$ is $A$-balanced and detects
\[
 U_\ell\otimes1-1\otimes U_\ell\ne0,
\]
while multiplication kills this element. It also annihilates every
bar boundary, proving that $B\widehat\otimes_A^{\mathbf L}B\to B$
is not a relative quasi-isomorphism.

There is a different, invertible comparison with a specified equivariant
domain. Let $C$ use the full new coefficients but only the old lattice,
and let $G$ be the added prime lattice. Then the actual new algebra is
$B=C\rtimes_{\rm rapid}G$. For complete $C$-bimodules with tempered
semilinear actions $\rho_g$, the maps
\[
 (mU_g)\otimes(nU_h)\longmapsto
       (m\otimes\rho_g(n))U_{g+h},\qquad
 (m\otimes n)U_k\longmapsto(mU_0)\otimes(nU_k)
\]
are inverse topological bimodule maps. Balancing over the actual gauge
unitaries proves the inverse identities, and tempered bounds absorb
the action in stronger rapid seminorms. They preserve the unit and
satisfy the associativity diagrams.

The comparison also holds for relative derived composition:
\[
 (M\rtimes G)\widehat\otimes_B^{\mathbf L}(N\rtimes G)
       \simeq (M\widehat\otimes_C^{\mathbf L}N)\rtimes G.
\]
Indeed the equivariant double bars remain resolutions after crossing;
their terms are free on each individual $B$ side, by
$(cU_g)\otimes v\mapsto(c\otimes\rho_gv)U_g$.
Thus their bounded-above complexes are relatively K-flat and compute
the asserted derived tensors. This repairs transport through the new
lattice on the declared equivariant domain. It still requires actual
$C$-bimodule lifts and new-generator actions on the geometric objects;
these data and their determinant descent are not supplied by induction.

These advances supply explicit local compositions, derived comparisons,
a relative determinant germ and controlled boundary metric data. Full
row (a) still requires their common geometric realization, the full mixed
diagonal and determinant comparison, and compatibility
with gluing, multiple-zero strata and the archimedean sector. The mixed
Riemann--Roch and global sign obligations of the other rows remain open.
